\chapter{Organ yang Bekerja}\label{ch:g7:organs-at-work}

Berlarilah cepat menaiki tiga lantai tangga lalu periksa keadaanmu:
\omterm{def:g1:my-body:parts}{tungkai} terbakar, \omterm{def:g4:heart-and-blood:heart}{jantung} berdentam, napas berpacu --- dan, anehnya,
\omterm{def:g1:the-five-senses:sense}{kulit} memerah dan hangat. Tidak seorang pun memerintahkan perubahan
itu; ia memerintahkan dirinya sendiri, dan dalam perbandingan yang
sempurna. Bab ini mengukur apa yang sesungguhnya dilahap sebuah organ
yang bekerja, dan bagaimana tubuh menata ulang pengantarannya pada
saat pekerjaannya bermula.

\section{Apa yang dilahap otot yang bekerja}

\begin{proposition}[Tagihan sang otot]\label{prop:g7:organs-at-work:bill}
\omterm{def:g3:bones-and-muscles:muscle}{Otot} yang bekerja membelanjakan sesuatu, dan pembelanjaannya dapat
diukur. Dibandingkan dengan \omterm{def:g3:bones-and-muscles:muscle}{otot} yang sama ketika beristirahat, \omterm{def:g3:bones-and-muscles:muscle}{otot}
yang bekerja keras:
\begin{itemize}
  \item mengambil \emph{\omterm{def:g7:what-is-respiration:respiration}{oksigen}} jauh lebih banyak dari \omterm{prop:g4:journey-of-food:absorb}{darah} yang
        melintasinya --- beberapa kali lipat timbaan istirahatnya;
  \item mengambil \emph{glukosa}\index{glukosa} jauh lebih banyak ---
        \omterm{def:g4:journey-of-food:digestion}{zat gizi} bergula yang siap pakai di dalam \omterm{prop:g4:journey-of-food:absorb}{darah}, bahan bakar
        sehari-harinya;
  \item melepaskan \emph{\omterm{def:g7:what-is-respiration:respiration}{karbon dioksida}} yang sepadan lebih banyak
        --- dan \emph{panas}.
\end{itemize}
Inilah pembakaran seluler pada
\cref{prop:g7:what-is-respiration:power}, yang dibaca dari alat ukur
sebuah organ: bahan bakar dan \omterm{def:g7:what-is-respiration:respiration}{oksigen} masuk, pekerjaan, gas buangan
dan kehangatan keluar.
\end{proposition}
\admitted

\begin{example}[Mengukur melintasi sebuah otot]\label{ex:g7:organs-at-work:measure}
Ahli faal membandingkan \omterm{prop:g4:journey-of-food:absorb}{darah} yang \emph{memasuki} sebuah \omterm{def:g3:bones-and-muscles:muscle}{otot} dengan
\omterm{prop:g4:journey-of-food:absorb}{darah} yang \emph{meninggalkannya}. Ketika beristirahat, \omterm{prop:g4:journey-of-food:absorb}{darah} yang
meninggalkannya sudah menyerahkan sebagian kecil \omterm{def:g7:what-is-respiration:respiration}{oksigen} dan
\omterm{prop:g7:organs-at-work:bill}{glukosanya}. Ketika bekerja keras, perbandingan yang sama
memperlihatkan \omterm{prop:g4:journey-of-food:absorb}{darah} yang meninggalkannya terkuras --- sebagian besar
\omterm{def:g7:what-is-respiration:respiration}{oksigen} yang dapat diantarkannya sudah diambil --- sementara \omterm{def:g7:what-is-respiration:respiration}{karbon
dioksidanya} melonjak. Organ yang sama, alat ukur yang sama: hanya
pekerjaannya yang berubah, dan tagihannya mengikutinya.
\end{example}

\begin{remark}[Kehangatannya adalah kuitansinya]\label{rem:g7:organs-at-work:warmth}
Panas tubuh yang bekerja bukanlah kebetulan: membakar bahan bakar
selalu menghasilkan kehangatan, dan \omterm{def:g3:bones-and-muscles:muscle}{otot} adalah ruang tungku tubuh.
Menggigil adalah kiat yang dijalankan dengan sengaja --- \omterm{def:g3:bones-and-muscles:muscle}{otot} yang
dipekerjakan tanpa tujuan lain selain kehangatan
(beban tetap pada \cref{ex:g3:balanced-diet:rest}, yang sejenak
dilipatduakan). \omterm{def:g1:the-five-senses:sense}{Kulit} yang memerah ketika berlari cepat adalah sistem
pendingin yang mengangkut panas tungkunya keluar --- urusan bagian
berikutnya.
\end{remark}

\section{Pengantaran sesuai permintaan}

\begin{proposition}[Tubuh menata ulang pasokannya]\label{prop:g7:organs-at-work:supply}
Pada saat pekerjaannya bermula, jalur pasokannya menata ulang diri
untuk membayar tagihannya (\cref{ex:g4:heart-and-blood:effort} sudah
mengumumkannya; inilah pengiriman yang lengkapnya):
\begin{enumerate}
  \item \emph{\omterm{def:g7:what-is-respiration:respiration}{pernapasan}} mendalam dan mencepat --- memuat ulang
        \omterm{prop:g4:journey-of-food:absorb}{darahnya} lebih cepat di \omterm{def:g7:how-animals-breathe:four}{paru-parunya}
        (\cref{ex:g4:breathing:running});
  \item \emph{\omterm{def:g4:heart-and-blood:heart}{jantung}} berdenyut lebih cepat dan lebih keras --- lebih
        banyak ronde per menit
        (\cref{prop:g4:heart-and-blood:pulse});
  \item \omterm{prop:g4:journey-of-food:absorb}{darahnya} \emph{dialihkan}: \omterm{def:g4:heart-and-blood:vessels}{pembuluh} yang melayani \omterm{def:g3:bones-and-muscles:muscle}{otot} yang
        bekerja (dan \omterm{def:g1:the-five-senses:sense}{kulit} yang mendinginkan) melebar, sementara
        \omterm{def:g4:heart-and-blood:vessels}{pembuluh} \omterm{def:g4:journey-of-food:tract}{saluran pencernaan} yang sedang beristirahat
        menyempit --- \omterm{prop:g4:journey-of-food:absorb}{darah} yang sama, yang dijalurkan ulang ke lantai
        yang sedang membelanjakan.
\end{enumerate}
Ketiganya mengikuti pekerjaannya secara otomatis, yang dikemudikan
oleh layanan \omterm{def:g5:brain-in-command:system}{otak} yang tidak diminta
(\cref{rem:g5:brain-in-command:more}).
\end{proposition}
\admitted

\begin{example}[Pengirimannya, yang dirasakan dari dalam]\label{ex:g7:organs-at-work:felt}
Pemeriksaan keadaan sesudah lari cepat menaiki tangga itu, yang
dijelaskan baris demi baris: napas yang berpacu --- pemuatan ulang;
\omterm{def:g4:heart-and-blood:heart}{jantung} yang berdentam --- ronde yang lebih cepat; \omterm{def:g1:my-body:parts}{tungkai} yang
terbakar --- tagihan \omterm{def:g3:bones-and-muscles:muscle}{ototnya} yang sedang dibayar sampai ke batasnya;
\omterm{def:g1:the-five-senses:sense}{kulit} hangat yang memerah --- panas tungkunya yang diangkut keluar ke
permukaan pendinginnya. Bahkan nyeri \omterm{def:g4:journey-of-food:tract}{lambung} yang klasik dan rasa
lesu sesudah santapan besar pun cocok: \omterm{def:g4:journey-of-food:digestion}{pencernaan} dan lari cepat
bersaing memperebutkan \omterm{prop:g4:journey-of-food:absorb}{darah} yang dijalurkan ulang, dan itulah
sebabnya \cref{met:g4:journey-of-food:habits} menasihatkan agar tidak
berlomba dengan perut yang penuh.
\end{example}

\begin{omfigure}
\begin{tikzpicture}
\begin{axis}[omaxis,
    width=9.6cm, height=5.8cm,
    xlabel={menit}, ylabel={denyut jantung per menit},
    xmin=0, xmax=13.6, ymin=40, ymax=200,
    xtick={0,2,4,6,8,10,12},
    ytick={60,100,140,180},
    ]
  \addplot[very thick, omThm, smooth] coordinates {
    (0,72) (1,74) (2,120) (3,150) (4,158) (5,160)
    (6,120) (7,100) (8,88) (9,80) (10,76) (11,74) (12,73)
  };
  \draw[dashed, omExo] (axis cs:2,40) -- (axis cs:2,180);
  \draw[dashed, omExo] (axis cs:5,40) -- (axis cs:5,180);
  \node[font=\footnotesize] at (axis cs:1.0,52) {istirahat};
  \node[font=\footnotesize] at (axis cs:3.5,52) {usaha};
  \node[font=\footnotesize] at (axis cs:8.6,52) {pemulihan};
\end{axis}
\end{tikzpicture}

{\small Sebuah \omterm{prop:g4:heart-and-blood:pulse}{denyut nadi} sepanjang tiga menit usaha yang keras:
tanjakannya bersama pekerjaannya, dataran tingginya dekat maksimumnya,
dan kembalinya yang perlahan --- \omterm{def:g7:organs-at-work:recovery}{pemulihan} --- sementara tubuhnya
melunasi rekeningnya.}
\end{omfigure}

\begin{definition}[Pemulihan]\label{def:g7:organs-at-work:recovery}
Ketika usahanya berhenti, \omterm{prop:g4:heart-and-blood:pulse}{denyut nadi} dan \omterm{def:g7:what-is-respiration:respiration}{pernapasannya} tidak: keduanya
turun kembali hanya secara bertahap --- yaitu tahap
\emph{pemulihan}\index{pemulihan}. Tubuhnya sedang melunasi utang:
memuat ulang cadangan \omterm{def:g7:what-is-respiration:respiration}{oksigennya}, membersihkan buangan usahanya,
melepaskan panas yang menumpuk. Makin bugar tubuhnya, makin cepat
pemulihannya --- dan itulah sebabnya waktu pemulihan itu sendiri
merupakan sebuah pengukuran kebugaran.
\end{definition}

\section{Latihan, dan batas mesinnya}

\begin{proposition}[Latihan menaikkan langit-langitnya]\label{prop:g7:organs-at-work:training}
Kalau dilatih secara teratur, seluruh rantai pasokannya menguat
(\cref{rem:g4:heart-and-blood:rest},
\cref{met:g4:breathing:care}): \omterm{def:g4:heart-and-blood:heart}{jantungnya} memindahkan lebih banyak
\omterm{prop:g4:journey-of-food:absorb}{darah} tiap denyut, \omterm{def:g3:bones-and-muscles:muscle}{otot} \omterm{def:g7:what-is-respiration:respiration}{pernapasannya} mendalam, \omterm{def:g3:bones-and-muscles:muscle}{ototnya} sendiri tumbuh
dan menyimpan lebih banyak bahan bakar, dan \omterm{def:g4:heart-and-blood:vessels}{pembuluh} halusnya berlipat
ganda. Hasilnya: usaha yang sama berharga \omterm{prop:g4:heart-and-blood:pulse}{denyut nadi} yang lebih
rendah bagi tubuh yang terlatih, dan langit-langitnya --- yaitu usaha
terberat yang dapat dibayarnya --- naik. Mesinnya dibangun ulang oleh
pemakaiannya sendiri.
\end{proposition}
\admitted

\begin{example}[Dua pelari, satu bukit]\label{ex:g7:organs-at-work:runners}
Yang tidak terlatih dan yang terlatih menaiki bukit yang sama. Yang
tidak terlatih: \omterm{prop:g4:heart-and-blood:pulse}{denyut nadi} $180$ di pertengahannya, terengah-engah,
terpaksa berjalan --- jalur pasokannya tidak dapat membayar
tagihannya, dan utangnya memanggil untuk berhenti. Yang terlatih:
\omterm{prop:g4:heart-and-blood:pulse}{denyut nadi} $150$ di puncaknya, masih dapat berbicara, pulih dalam
tiga menit. Bukit yang sama, tagihan yang sama --- langit-langit yang
berbeda, yang dibangun berbulan-bulan lewat bukit semacam itu juga
yang didaki pelan-pelan.
\end{example}

\begin{method}[Membaca mesinmu sendiri]\label{met:g7:organs-at-work:read}
Sebuah percobaan atas diri sendiri dengan sebuah jam
(\cref{met:g4:heart-and-blood:pulse} menyediakan teknik \omterm{prop:g4:heart-and-blood:pulse}{denyut
nadinya}):
\begin{enumerate}
  \item catatlah \omterm{prop:g4:heart-and-blood:pulse}{denyut nadi} istirahatmu, sambil duduk dan tenang;
  \item lakukan tiga menit usaha yang mantap --- naik-turun bangku,
        lompat tali;
  \item catatlah \omterm{prop:g4:heart-and-blood:pulse}{denyut nadinya} seketika, lalu tiap menit sampai ia
        kembali ke keadaan istirahat: itulah waktu \omterm{def:g7:organs-at-work:recovery}{pemulihannya};
  \item ulangi setiap minggu sepanjang satu triwulan olahraga yang
        teratur: amati \omterm{prop:g4:heart-and-blood:pulse}{denyut nadi} istirahat, \omterm{prop:g4:heart-and-blood:pulse}{denyut nadi} usaha dan
        waktu \omterm{def:g7:organs-at-work:recovery}{pemulihannya} semuanya menurun perlahan --- pembangunan
        ulang pada \cref{prop:g7:organs-at-work:training}, dalam
        datamu sendiri.
\end{enumerate}
\end{method}

\section{Latihan}

\begin{exercise}[$\star$]\label{exo:g7:organs-at-work:1}
Daftarkan apa yang lebih banyak diambil sebuah \omterm{def:g3:bones-and-muscles:muscle}{otot} yang bekerja, dan
apa yang lebih banyak dilepaskannya.
\end{exercise}

\begin{exercise}[$\star$]\label{exo:g7:organs-at-work:2}
Apa itu \omterm{prop:g7:organs-at-work:bill}{glukosa}, dan dari mana \omterm{def:g3:bones-and-muscles:muscle}{otot} yang bekerja memperolehnya?
\end{exercise}

\begin{exercise}[$\star$]\label{exo:g7:organs-at-work:3}
Lukiskan perbandingan \omterm{prop:g4:journey-of-food:absorb}{darah} yang memasuki melawan yang meninggalkan,
dan apa yang diperlihatkannya ketika beristirahat melawan ketika
berusaha.
\end{exercise}

\begin{exercise}[$\star$]\label{exo:g7:organs-at-work:4}
Sebutkan ketiga penataan ulang pasokan pada
\cref{prop:g7:organs-at-work:supply}.
\end{exercise}

\begin{exercise}[$\star$]\label{exo:g7:organs-at-work:5}
Mengapa pekerjaan yang berat menghangatkanmu --- dan menggigil itu
apa, dalam istilah bab ini?
\end{exercise}

\begin{exercise}[$\star$]\label{exo:g7:organs-at-work:6}
Apa itu \omterm{def:g7:organs-at-work:recovery}{pemulihan}, dan utang apa saja yang sedang dilunasi tubuhnya
selama \omterm{def:g7:organs-at-work:recovery}{pemulihan} itu?
\end{exercise}

\begin{exercise}[$\star\star$]\label{exo:g7:organs-at-work:7}
Bacalah grafik \omterm{prop:g4:heart-and-blood:pulse}{denyut nadinya}: nilai istirahatnya, nilai dataran
tingginya, dan kira-kira berapa lama \omterm{def:g7:organs-at-work:recovery}{pemulihannya}. Apa yang akan
diubah latihan pada masing-masingnya?
\end{exercise}

\begin{exercise}[$\star\star$]\label{exo:g7:organs-at-work:8}
Jelaskan rasa lesu sesudah makan siang dan nyeri \omterm{def:g4:journey-of-food:tract}{lambung} si pelari
cepat dengan aturan penjaluran ulangnya.
\end{exercise}

\begin{exercise}[$\star\star$]\label{exo:g7:organs-at-work:9}
Mengapa \omterm{def:g1:the-five-senses:sense}{kulit} yang memerah ketika berusaha merupakan perubahan
\emph{pengantaran}, bukan sebuah kerusakan? Ia sedang memecahkan
masalah siapa?
\end{exercise}

\begin{exercise}[$\star\star$]\label{exo:g7:organs-at-work:10}
Daftarkan tiga pembangunan ulang yang dengannya latihan menaikkan
langit-langitnya, dan tanda sehari-hari dari masing-masingnya.
\end{exercise}

\begin{exercise}[$\star\star$]\label{exo:g7:organs-at-work:11}
Dua murid menjalankan \cref{met:g7:organs-at-work:read}: \omterm{prop:g4:heart-and-blood:pulse}{denyut nadi}
istirahatnya $68$ dan $86$; \omterm{def:g7:organs-at-work:recovery}{pemulihannya} tiga dan delapan menit. Mesin
mana yang terlatih, dan apa tepatnya yang dijamin masing-masing kedua
bilangan itu?
\end{exercise}

\begin{exercise}[$\star\star\star$]\label{exo:g7:organs-at-work:12}
Hubungkan tiga bab dalam satu alinea: tagihan sang \omterm{def:g3:bones-and-muscles:muscle}{otot}
(\cref{prop:g7:organs-at-work:bill}), tujuan \omterm{def:g7:what-is-respiration:respiration}{respirasi}
(\cref{prop:g7:what-is-respiration:power}), dan layanan pengantaran
\omterm{prop:g4:journey-of-food:absorb}{darahnya}
(\cref{prop:g4:heart-and-blood:carries}) --- sebagai satu penjelasan
tentang satu kali lari cepat.
\end{exercise}

\section{Soal: Kajian Kebugaran Sekolah}

\begin{problem}[{Soal akhir pekan --- setriwulan data denyut nadi,
yang dibaca seperti seorang ahli faal}]\label{pb:g7:organs-at-work:1}
Kelasnya menjalankan \cref{met:g7:organs-at-work:read} setiap hari
Senin selama satu triwulan: \omterm{prop:g4:heart-and-blood:pulse}{denyut nadi} istirahat, \omterm{prop:g4:heart-and-blood:pulse}{denyut nadi} sesudah
tiga menit naik-turun bangku, dan waktu \omterm{def:g7:organs-at-work:recovery}{pemulihan}. Dua catatan yang
disamarkan namanya: murid A --- minggu 1: istirahat $84$, usaha $168$,
\omterm{def:g7:organs-at-work:recovery}{pemulihan} $9$ menit; minggu 12: istirahat $72$, usaha $150$, \omterm{def:g7:organs-at-work:recovery}{pemulihan}
$5$ menit. Murid B --- minggu 1: istirahat $70$, usaha $150$,
\omterm{def:g7:organs-at-work:recovery}{pemulihan} $4$ menit; minggu 12: istirahat $69$, usaha $148$, \omterm{def:g7:organs-at-work:recovery}{pemulihan}
$4$ menit.

\medskip\noindent\textbf{Bagian I --- Di dalam satu sesi.}
\begin{enumerate}
  \item Selama naik-turun bangkunya, daftarkan apa yang sedang terjadi
        pada timbaan \omterm{def:g7:what-is-respiration:respiration}{oksigen} dan \omterm{prop:g7:organs-at-work:bill}{glukosa} sebuah \omterm{def:g3:bones-and-muscles:muscle}{otot} paha, dan pada
        keluaran \omterm{def:g7:what-is-respiration:respiration}{karbon dioksida} serta panasnya.
  \item Jelaskan mengapa \omterm{def:g7:what-is-respiration:respiration}{pernapasan} dan \omterm{prop:g4:heart-and-blood:pulse}{denyut nadinya} menanjak
        \emph{bersama-sama} selama menit-menit usahanya.
  \item Seorang murid merasa hangat dan memerah menjelang menit
        ketiga. Berikan penjelasan tungku-dan-pendinginannya.
  \item Mengapa \omterm{prop:g4:heart-and-blood:pulse}{denyut nadinya} tidak turun ke keadaan istirahat pada
        saat naik-turun bangkunya berhenti?
\end{enumerate}

\medskip\noindent\textbf{Bagian II --- Membaca triwulannya.}
\begin{enumerate}[resume]
  \item Lukiskan perubahan A sepanjang triwulannya, lalu sebutkan
        pembangunan ulang pada
        \cref{prop:g7:organs-at-work:training} di balik tiap
        bilangannya.
  \item Bilangan B nyaris tidak bergerak. Berikan kedua pembacaan yang
        jujur --- dan fakta tambahan apa tentang kehidupan B di luar
        kelas yang akan memutuskan di antara keduanya.
  \item Pada minggu ke-5, seorang murid mengukur tepat sesudah makan
        siang yang tergesa-gesa lalu memperoleh bilangan yang ganjil.
        Aturan \cref{prop:g7:organs-at-work:supply} yang mana yang
        mengganggu, dan apa yang harus dikatakan tata caranya tentang
        waktu pengukuran?
  \item Mengapa tiap murid harus selalu mengukur dengan cara yang sama
        --- usaha yang sama, pengaturan waktu yang sama? (Seorang
        kawan lama di antara metodenya; sebutkan namanya.)
\end{enumerate}

\medskip\noindent\textbf{Bagian III --- Surat sang ahli
faal.}
\begin{enumerate}[resume]
  \item Tuliskan rangkuman dua kalimat tentang triwulan A untuk papan
        pengumuman kelas, dengan memakai \emph{langit-langit} dan
        \emph{\omterm{def:g7:organs-at-work:recovery}{pemulihan}}.
  \item Seorang orang tua bertanya apakah \omterm{prop:g4:heart-and-blood:pulse}{denyut nadi} istirahat yang
        rendah itu mengkhawatirkan. Jawablah dengan pengendara sepeda
        pada \cref{exo:g4:heart-and-blood:11}.
  \item Kelasnya menginginkan satu bilangan untuk dilacak triwulan
        berikutnya. Berikan alasan bagi waktu \omterm{def:g7:organs-at-work:recovery}{pemulihan} alih-alih
        \omterm{prop:g4:heart-and-blood:pulse}{denyut nadi} usaha, dengan memakai
        \cref{def:g7:organs-at-work:recovery}.
  \item Tutuplah dengan pelajaran moral kajiannya dalam satu kalimat:
        apa yang membangun ulang mesin A?
\end{enumerate}
\end{problem}
